% ==========================================================
% titles/title-08.tex -- Title VIII: Borders and Membership
% ==========================================================

\ConstTitle{Borders and Membership}{title:borders}

\Article{Open Physical Presence and Universal Rights}{art:open-presence}

\ArtSection{Open Physical Presence}{art:open-presence:s-open-physical-presence}
Any person may enter Commonwealth territory and remain. Presence is not criminalized. There is no detention for movement. The coercive border enforcement characteristic of prior systems, which sorted persons into radically different life outcomes based on the morally arbitrary circumstance of birthplace, is abolished.

\ArtSection{Universal Rights Upon Arrival}{art:open-presence:s-universal-rights-upon}
Every person present on Commonwealth territory is immediately entitled, from the moment of arrival and without application, means-testing, or demonstration of future intent to remain, to the material rights guaranteed under Article \ref{art:matsuff}, \ref{art:matsuff:s-guaranteed-material-rights}, together with:
\begin{itemize}
  \item Full legal protections and due process rights
  \item Freedom from arbitrary detention or removal
\end{itemize}

\ArtSection{Commons-Allocated Housing}{art:open-presence:s-commons-allocated-housing}
Commons-Allocated Housing is the primary mechanism through which all persons on Commonwealth territory access stable housing. The \gls{gls:housing-implementation-council} maintains and administers a Commons housing stock sufficient to meet demand, governed by the following principles:

Allocation is determined by published objective criteria including household size, accessibility requirements, proximity to work or care responsibilities, and expressed preferences. No criterion related to civic integration status, race, gender, sexual orientation, religion, national origin, or any characteristic protected under Article \ref{art:equality} may be considered in allocation decisions. Where multiple applicants meet criteria equally, allocation is determined by sortition.

All allocation decisions are logged in real time through the \gls{gls:transparency-engine}. Persons denied a specific allocation receive a written determination stating the criteria applied. Appeals are heard by the \gls{gls:housing-appeals-panel}, a standard body of nine members serving three-year terms, with no membership overlap with the \gls{gls:housing-implementation-council}. \gls{gls:constitutional-court} review is available on constitutional grounds.

The \gls{gls:housing-implementation-council} publishes annual reports through the \gls{gls:transparency-engine} on housing stock levels, allocation waiting times, demographic patterns in allocation outcomes, and progress toward eliminating any housing gap. Systematic disparities in allocation outcomes along any protected characteristic trigger mandatory review by the \gls{gls:minority-impact-assessment-body} under Article \ref{art:local-democracy}, \ref{art:local-democracy:s-minority-impact-assessment}.

\ArtSection{Persons Fleeing Persecution or Crisis}{art:open-presence:s-persons-fleeing-persecution}
Persons fleeing persecution, violence, or acute crisis have an explicit constitutional guarantee that their \ref{art:open-presence:s-universal-rights-upon} rights continue for the duration of their danger regardless of duration of presence. This guarantee requires no application, determination, or bureaucratic process. The Commonwealth acknowledges that many persons fleeing danger are fleeing conditions that wealthy nations helped create, and that this history generates particular obligations of welcome. The prohibition on arbitrary detention or removal applies with full force to all persons on Commonwealth territory; no person may be removed on the basis of their persecution or crisis status, or lack thereof.

\ArtSection{The Anti-Domination Basis}{art:open-presence:s-anti-domination-basis}
These rights rest on the same anti-domination rationale as Article \ref{art:matsuff}: material desperation is a lever of control, and the Commonwealth removes that lever entirely. Here, that guarantee begins for every person on Commonwealth territory from the moment of their arrival, without regard to duration of presence or immigration status.

\ArtSection{Stable Housing Transition Obligation}{art:open-presence:s-stable-housing-transition}
The \gls{gls:housing-implementation-council} maintains a public vacancy register through the \gls{gls:transparency-engine} and a private register of Residential Tenures abandoned for over three years. The three-year threshold period is public and subject to \gls{gls:sortition-legislature} review. Where the \gls{gls:housing-implementation-council} determines that a housing gap exists -- defined as any person on Commonwealth territory lacking access to stable housing or transitional housing meeting the standards of this Article -- that cannot be fulfilled by the public vacancy register, all abandoned Residential Tenures are purchased by the Commons and the relevant properties are transferred to Commons housing stock.

Multi-family residences with private sleeping quarters and shared facilities are recognized as a transitional housing form meeting a dignity standard between emergency shelter and fully stable housing. They may be provided as transitional accommodation where fully stable housing cannot be immediately supplied, subject to the following conditions:
\begin{itemize}
  \item Residents retain all rights under \ref{art:open-presence:s-universal-rights-upon}
  \item The \gls{gls:justice-standards-council} sets minimum dignity standards for transitional multi-family residences
  \item The \gls{gls:housing-implementation-council} is constitutionally obligated to provide each resident in transitional multi-family accommodation with fully stable housing within three years of their entry into that accommodation, with annual public reporting on progress
\end{itemize}

Where stable housing cannot be immediately provided due to supply constraints at the time of ratification or thereafter, emergency shelter of equivalent dignity is provided as a transitional measure. Emergency shelter is never a permanent substitute for stable housing. The \gls{gls:housing-implementation-council} (Article \ref{art:implementation-councils}) is constitutionally obligated to eliminate any gap between emergency shelter provision and stable housing provision within five years of ratification, with annual public reporting on progress through the \gls{gls:transparency-engine}. New gaps created as a result of immigration to the Commonwealth must be addressed by the \gls{gls:housing-implementation-council} within two years of the gap being identified and reported on. The \gls{gls:sortition-legislature} may extend either transition period once only, by supermajority vote, if supply constraints genuinely prevent achievement within five years (ratification) or two years (reported new gap). No further extension is permitted. All persons in emergency shelter transition accommodations retain full rights under \ref{art:open-presence:s-universal-rights-upon} throughout.

\Article{Residency, Civic Integration, and Full Membership}{art:civic-integration}

\ArtSection{Automatic Residency}{art:civic-integration:s-automatic-residency}
Any person present on Commonwealth territory for more than ninety continuous days automatically receives residency status. Residency is not applied for and not subject to bureaucratic discretion. It is automatic and immediate upon meeting the time condition.

\ArtSection{Residency Rights}{art:civic-integration:s-residency-rights}
Residency status confers full access to the Material Sufficiency Floor, the right to work and join cooperatives, access to all Commons services, and absolute freedom from deportation or forced removal. Residency retains full legal protections and due process rights already provided to visitors.

\ArtSection{The Civic Integration Index}{art:civic-integration:s-civic-integration-index}
Full political membership may be achieved through the Civic Integration Index. Its dimensions are:
\begin{itemize}
  \item Active cooperative participation as a worker-member, verified through cooperative membership records held by the relevant cooperative and confirmed by the \gls{gls:electoral-commons} under Article \ref{art:sortition-design}, \ref{art:sortition-design:s-electoral-commons}
  \item Completion of the constitutional principles education program, offered free at any pace in all threshold languages, verified through the education Commons institution upon completion
\end{itemize}

These dimensions are chosen because they are verifiable through existing institutional records without requiring behavioral monitoring of residents. The Commonwealth explicitly rejects participation tracking systems that would require surveillance of residents' civic behavior, recognizing that such systems are inconsistent with Article \ref{art:internal-security}, \ref{art:internal-security:s-no-mass-surveillance}'s prohibition on mass surveillance and with the anti-domination principle. If future technical architecture enables privacy-respecting civic engagement tracking that is genuinely consistent with these principles, the \gls{gls:sortition-legislature} may add additional Index dimensions by supermajority vote subject to \gls{gls:constitutional-court} review confirming consistency with Article \ref{art:transparency} and Article \ref{art:civic-integration}.

Worker-membership in a cooperative, including full ownership rights and democratic governance participation within that cooperative, is available to all residents of the Commonwealth regardless of civic integration status. Cooperative ownership rights are governed by Article \ref{art:worker-coop-sector} and are distinct from political membership rights conferred by full civic integration under this Article. Cooperative worker-membership by residents counts toward the Civic Integration Index under this Section.

\ArtSection{Flexible Achievement}{art:civic-integration:s-flexible-achievement}
The Index may be achieved at any pace. A minimum residency floor of one year applies. Someone who arrives and engages immediately in cooperative work may achieve full membership in approximately eighteen months upon education program completion. Someone dealing with trauma, language acquisition, disability, caregiving responsibilities, or other legitimate life circumstances may take longer without penalty, and remains on the automatic membership pathway of \ref{art:civic-integration:s-maximum-time-ceiling} throughout. The Commonwealth recognizes community connection -- participation in mutual aid networks, local Commons governance, and neighborhood life -- as a valued dimension of civic culture that the Index does not and cannot measure. The Commons infrastructure actively supports these forms of participation because they strengthen the democratic culture this Constitution is designed to sustain, not because they are conditions of membership.

\ArtSection{Maximum Time Ceiling}{art:civic-integration:s-maximum-time-ceiling}
Any person who has maintained residency for four continuous years is automatically granted full civic membership regardless of Index score, provided they have not been found by the \gls{gls:constitutional-court} to have actively and deliberately worked against the foundational principles during that period. Passive non-participation does not disqualify.

\ArtSection{Automatic Conferral}{art:civic-integration:s-automatic-conferral}
Upon meeting either the Index threshold or the four-year ceiling, full civic membership is conferred automatically. It is not subject to bureaucratic discretion or rejection. Membership confers the right to be selected for governance positions via sortition and the right to vote in all national direct votes, referenda, and civic initiative processes under this Constitution. Residents who have not yet attained full civic membership retain the full rights of Title I and Title VIII regardless, but do not vote in national processes until membership is conferred.

\ArtSection{No Cultural Assimilation Requirement}{art:civic-integration:s-no-cultural-assimilation}
Full membership does not require language fluency beyond functional communication, cultural conformity, or value alignment beyond genuine commitment to the foundational principles of this Constitution -- including the equality provisions of Article \ref{art:equality}.

\ArtSection{Foreign Funding Prohibition}{art:civic-integration:s-foreign-funding-prohibition}
All political organizations, media outlets, civil society groups, and advocacy organizations must publicly disclose all funding sources in real time through the \gls{gls:transparency-engine}. Foreign funding of domestic political activity is prohibited and criminally enforceable.